\scanpage{55}
% typed verso
% @top centre, pagination
\mspage{- 130 -}
% @top right, running head
\hfill IX

% @01–02, opening
Elle implique bien que si $\varphi^t_{\xi_{\unsure{o|0}}}$ est symétrique et positive, il en est de même de $\varphi^t_\xi$.

% @03–11, nondegeneracy
Enfin, le fait que $\varphi^t_{\xi_o}$ soit non dégénérée équivaut évidemment à l'assertion que l'accouplement (10.4.1) est séparant (compte tenu que $\xi_{ot}:M\otimes_{\mathbb{Z}}\mathbb{Q}\handwritten{\longrightarrow}M'\otimes_{\mathbb{Z}}\mathbb{Q}$ est un isomorphisme, $\widetilde{\xi}_o$ étant une isogénie), et est donc également indépendant de la polarisation choisie. On peut d'ailleurs noter que pour vérifier que (10.4.1) est séparant, il suffit de vérifier la même assertion pour l'accouplement de monodromie $u_\ell$ (9.1.2) pour \uline{un} nombre premier $\ell$. Prenant $\ell\ne p$, on peut utiliser la description 9.2 de $u_\ell$ ; le fait que $u_\ell$ soit séparant à gauche (donc des deux côtés, pour une raison de rangs) est alors trivial.

% @12–13, section title
11. \uline{Composantes connexes du modèle de Néron} : \uline{relation de dualité}, \uline{comportement assymptotique}\unsure{.}

% @14, opening of section
11.0. Le but du présent paragraphe est l'étude du $k$-groupe fini étale
% @15, component group
\[
 \Phi_o=A_o/A_o^o\tag*{(11.0.1)}
\]
% @16, introducing the primary component
déjà considéré dans 1. Pour tout nombre premier $\ell$, nous désignons par
% @17, primary component
\[
 \Phi_o(\ell)
\]
% @18–19, decomposition
la composante $\ell$-primaire de ce dernier, de sorte que la connaissance de $\Phi_o$ équivaut à celle des $\Phi_o(\unsure{1|l})$, grâce à la formule
% @20, decomposition formula
\[
 \Phi_o=\handwritten{\coprod}_{\ell}\Phi_o(\ell)\quad .\tag*{(11.0.2)}
\]
% @lower strip, pencil foliation
\margin{\handwritten{448}}
% @lower centre, printed label
\mspage{448}
% archive sheet number: 54

